\subsection{The Same Vector on Two Different Grids}

\begin{examplebox}
Let the ordinary Cartesian basis be
\[
\mathcal E=(\mathbf e_1,\mathbf e_2),
\qquad
\mathbf e_1=(1,0),
\qquad
\mathbf e_2=(0,1),
\]
and choose the slanted basis
\[
\mathcal B=(\mathbf b_1,\mathbf b_2),
\qquad
\mathbf b_1=(2,1),
\qquad
\mathbf b_2=(1,2).
\]
Consider the geometric vector
\[
\mathbf u=(5,4)
\]
in Cartesian coordinates. On the rectangular grid,
\[
[\mathbf u]_{\mathcal E}=(5,4).
\]
On the slanted grid, we seek $c_1,c_2$ such that
\[
\mathbf u=c_1\mathbf b_1+c_2\mathbf b_2.
\]
The two pictures show the same arrow. Only the grid used to read its coordinates
has changed.

\begin{center}
\begin{tikzpicture}[scale=0.72,>=stealth]
    \begin{scope}[xshift=0cm]
        \draw[gray!35,step=1] (-0.3,-0.3) grid (5.7,4.7);
        \draw[axis,->] (-0.3,0) -- (5.9,0) node[right] {$\mathbf e_1$};
        \draw[axis,->] (0,-0.3) -- (0,4.9) node[above] {$\mathbf e_2$};
        \coordinate (OL) at (0,0);
        \coordinate (UL) at (5,4);
        \draw[helper,dashed] (UL) -- (5,0);
        \draw[helper,dashed] (UL) -- (0,4);
        \draw[very thick,->] (OL) -- (UL)
            node[midway,above left] {$\mathbf u$};
        \node[below] at (5,0) {$5$};
        \node[left] at (0,4) {$4$};
        \node[align=center] at (2.7,-1.0)
            {rectangular grid\\$[\mathbf u]_{\mathcal E}=(5,4)$};
    \end{scope}

    \begin{scope}[xshift=10.2cm]
        \coordinate (OR) at (0,0);
        \coordinate (B1) at (2,1);
        \coordinate (B2) at (1,2);
        \coordinate (TwoB1) at (4,2);
        \coordinate (UR) at (5,4);

        \foreach \i in {-1,0,1,2,3} {
            \draw[gray!45,dashed]
                ({-1+2*\i},{-2+\i}) -- ({3+2*\i},{6+\i});
        }
        \foreach \j in {-1,0,1,2,3} {
            \draw[gray!45,dashed]
                ({-2+\j},{-1+2*\j}) -- ({6+\j},{3+2*\j});
        }

        \draw[blue!70!black,line width=1pt,->] (OR) -- (B1)
            node[midway,below right] {$\mathbf b_1$};
        \draw[orange!85!black,line width=1pt,->] (OR) -- (B2)
            node[midway,left] {$\mathbf b_2$};
        \draw[blue!70!black,line width=1pt,->] (OR) -- (TwoB1)
            node[midway,below right] {$2\mathbf b_1$};
        \draw[orange!85!black,line width=1pt,->] (TwoB1) -- (UR)
            node[midway,right] {$\mathbf b_2$};
        \draw[very thick,->] (OR) -- (UR)
            node[midway,above left] {$\mathbf u$};
        \draw[orange!85!black,dashed] (UR) -- (TwoB1);
        \draw[blue!70!black,dashed] (UR) -- (B2);
        \node[align=center] at (2.8,-1.0)
            {slanted grid\\$[\mathbf u]_{\mathcal B}=(2,1)$};
    \end{scope}
\end{tikzpicture}
\end{center}

On the rectangular grid the coordinates are read along perpendicular axes. On
the slanted grid the lines through the endpoint are drawn parallel to the other
basis vector, producing an oblique parallelogram rather than a rectangle.
\end{examplebox}
